
{\actuality} Современные интернет-сети становятся все более сложными, а также характеризуются высокой динамичностью изменений, что предъявляет высокие требования к методам их анализа и оценке производительности в режиме реального времени или близким к таковым по скорости. Эффективность работы сети напрямую зависит от множества факторов, включая нагрузку на каналы передачи данных, количество активных пользователей, процесс маршрутизации, а также используемые типы устройств. Все это влияет на качество обслуживания (QoS - Quality-of-Service) в условиях ограниченного вычислительного ресурса. С увеличением объема данных и числа подключенных устройств требуется разработка новых подходов для мониторинга и анализа работы сетей в реальном времени.

Однако в условиях ограниченного вычислительного ресурса, а также необходимости обеспечения быстродействия и масштабируемости, использование сложных моделей машинного обучения, таких как глубокие нейронные сети, оказывается не всегда оправданным. Эти модели требуют значительных вычислительных мощностей и памяти, что усложняет процесс внедрения в реальных системах с ограниченными ресурсами. Многие современные исследования указывают на необходимость в поиске альтернатив для задач сетевого мониторинга, в частности при применении устройств периферии сети (edge device) \cite{Xu2019LightweightNet, Wang2024Lightweight, Liu2022Analysis}. Также стоит подчеркнуть, что различные ресурсоемкие методы (в частности нейросетевые подходы) требуют больших серверных мощностей, а значит необходима передача данных с точки сбора до данного сервера, что может не проходить по ограничениям по времени, а также нарушать требования безопасности при транспортировке чувствительных данных. В связи с перечисленными проблемами требуется разработка и применение легковесных алгоритмов, которые позволяют на должном уровне точности анализировать производительность интернет-сетей при минимальных затратах ресурсов и таким образом не теряя свою эффективность с ростом сетевой инфраструктуры и увеличением количества пользователей.


%\textbf{Актуальность работы}

%\textbf{Степень разработанности темы}
%(Есть множество публикаций по теме, но в них идет фокус на использовании DNN, что значительно отличается от вектора данной работы. )

{\aim} данной работы является исследование и разработка методологии анализа записей сетевого трафика, методов прогнозирования времени отклика узлов сети для применения на периферийных сетевых устройствах с высокой вычислительной нагрузкой.

%данной работы является исследование и разработка методов обработки и анализа интернет-трафика и данных интернет-сетей, в частности анализ доступности соединения, с упором на использование легковесных методов и алгоритмов, в частности относящихся к технологиям TinyML.


Для достижения данной цели были поставлены следующие {\tasks}:
\begin{enumerate}[beginpenalty=10000] 
	\item Исследовать записи сетевого трафика из открытых источников,  провести сравнительный анализ методов \phantom{ } генерации \phantom{ } и преобразования признаков, провести разработку моделей классификаторов трафика с учётом ограничений вычислительной сложности, сформировать признаковое пространство для повышения метрик классификации трафика.
	\item Провести исследование распределенных систем измерений параметров интернет сетей, разработать методы предобработки и калибровки получаемых показателей.
	\item Разработать методы анализа и прогнозирования времени отклика при учёте применения периферийных сетевых устройств с высокой вычислительной нагрузкой. 
	%\item Провести исследование современных решений задачи анализа интернет-сетей, в частотности в условиях ограничений на вычислительные мощности. 
	%\item Исследовать данные о трафике из открытых источников и провести апробацию и сравнение методов анализа соединений, применяемых с целью обеспечения необходимого уровня QoS.
	%\item Провести сбор и исследование данных о состоянии сетевой инфраструктуры. Разработать методы предобработки полученных данных.
	%\item Разработать методы анализа и оценки состояния компьютерных сетей, в том числе в режиме реального времени. При ведении разработки учитывать малую мощность применяемых устройств: отказ от использования глубоких нейронных сетей (DNN) и высоконагруженных ансамблевых методов машинного обучения.  
	%\item \ldots
\end{enumerate}


%{\novelty}
%\begin{enumerate}[beginpenalty=10000] 
%	\item Предложен комплексный набор методов предобработки и калибровки получаемых показателей \ldots
%	\item Разработана система оценки загруженности сети, не использующая прямых измерений. \ldots%
%	\item Было выполнено оригинальное исследование на реальных данных работы сети на территории РФ, формирующее актуальную картину работы сетей\ldots
%	\item ...
%\end{enumerate}

{\defpositions}
\begin{enumerate}[beginpenalty=10000]
	\item Проведенный разведочный анализ, исследование и разработка методологии анализа данных сетевого домена позволили выявить ковариационный сдвиг в ряде датасетов и улучшить модели классификации (увеличение accuracy не менее чем на 0,2), а также выявить оптимальные методы предобработки.
	
	\item Анализ и исследование особенностей программно-аппаратных платформ системы RIPE Atlas привел к разработке методов предобработки и калибровки получаемых данных, которые позволяют компенсировать временную задержку у 63\% устройств, обеспечивая их функциональную эквивалентность 37\% эталонных устройств и совместимость с едиными аналитическими методами.
	
	\item Разработанные и исследованные алгоритмы динамической кластеризации измерительных зондов по показателям времени отклика на территории РФ с применением методов машинного обучения позволили выделить 32\% аномальных узлов, фильтрация которых привела к снижению ошибки (RMSE) прогнозирования времени приема-передачи до 50\%.
	
	\item Разработанные и исследованные алгоритмы прогнозирования времени приема-передачи (round-trip-time, RTT) в детектируемых кластерах измерительных зондов позволили использовать на 3 порядка меньше операций с плавающей запятой (FLOPs) при достижении сопоставимой точности в сравнении с передовыми нейросетевыми подходами.
	
\end{enumerate}

%\textbf{Практическая значимость/Теоретическая значимость}


{\reliability} \phantom{ } полученных результатов \phantom{ } обеспечивается использованием методов статистического анализа, машинного обучения, математическим моделированием, совпадением результатов исследования с экспериментальными данными, а также непосредственным участием автора в получении исходных данных и проведении экспериментов.


\textbf{Апробация работы.} Основные результаты работы докладывались на
следующих научных конференциях:
	\begin{enumerate}[beginpenalty=10000]
		\item VI международная конференция
		«Информационные технологии и технические средства управления» (ICCT-2022), Институт проблем
		управлений им. В.А.Трапезникова РАН совместно с Астраханским государственным
		техническим университетом, 2022 г.
		\item X Международная конференция  «Инжиниринг  \& \\ Телекоммуникации − En\&T-2023» МФТИ, Долгопрудный
		\item 65-й Всероссийская научная конференция
		МФТИ в честь 115-летия Л.Д. Ландау 2023, МФТИ, Долгопрудный
		\item XXVI Международная конференция
		«Цифровая обработка сигналов и ее применение — DSPA-2024», Институт проблем управления им. В.А.Трапезникова РАН, Москва, Россия.
		\item XI Международная конференция  «Инжиниринг \& \\ Телекоммуникации — En\&T-2024» МФТИ, Москва.
	\end{enumerate}
	
{\publications} Основные результаты по теме диссертации изложены в 6 печатных работах, 1 из которых издана в журналах, рекомендованных ВАК, 2 –
в рецензируемых изданиях, входящих в базу данных Scopus, 3 - в тезисах докладов. Получено 1 свидетельство о регистрации программы для ЭВМ.




